\documentclass[11pt]{article}
\usepackage[margin=1in]{geometry}
\usepackage{amsmath,amssymb}
\usepackage{booktabs}
\usepackage{hyperref}
\usepackage{siunitx}

\title{Independent Replication Report --- OSTI 3374709\\
\large Verification of an energy-conserving semi-implicit ES-PIC scheme}
\author{Independent replication (OSTI-100, rank 4, accelerator\_plasma)}
\date{}

\begin{document}
\maketitle

\section*{Paper}
R.\,M.~Hedlof, D.\,C.~Barnes, R.\,E.~Groenewald, \emph{et al.},
``Verification of an energy-conserving semi-implicit electrostatic
particle-in-cell scheme for modeling high-density plasma at scale,''
\textbf{Phys.\ Plasmas 33, 053902 (2026)}.
DOI: \href{https://doi.org/10.1063/5.0315721}{10.1063/5.0315721};
OSTI: 3374709; LBNL eScholarship item 8xt682g7; CC-BY 4.0.

\section{Paper summary}
Momentum-conserving explicit electrostatic PIC (ECPIC) must resolve the
electron plasma period ($\omega_{pe}\Delta t \lesssim 2$) for stability,
which is prohibitively expensive for high-density plasmas. The authors
present a \emph{semi-implicit} ES-PIC (SIPIC) in which the particle mass
is replaced by an SPD operator that suppresses only the Langmuir
(plasma-wave) response. In the implemented form (their Eqs.~9--10,
barycentric shape functions) this is exactly equivalent to solving
Poisson's equation with a density-dependent effective dielectric
\[
\varepsilon_{\text{eff}} = \varepsilon_0\!\left(1 + \tfrac{C_{\text{SI}}\,\omega_{pe}^{2}\,\Delta t^{2}}{4}\right),
\]
so that the electron plasma mode is down-shifted (their Eq.~16) to
\[
\omega_{pe,\text{SI}} = \frac{\omega_{pe}}{\sqrt{1 + C_{\text{SI}}\,\omega_{pe}^{2}\,\Delta t^{2}/4}}.
\]
This makes the scheme stable at $\Delta t$ far beyond $1/\omega_{pe}$ while
remaining second-order accurate for lower-frequency modes and (per
\S III.B) leaving cyclotron / hybrid physics essentially unaffected.
\S III demonstrates this against analytic dispersion (\S III.A
Bohm--Gross), upper/lower hybrid modes (\S III.B), Landau damping
(\S III.C), and energy conservation, using WarpX and Aleph code-to-code
comparison.

\section{Claims table}
\begin{center}\small
\begin{tabular}{@{}clccc@{}}
\toprule
ID & Claim & Type & Testable & Tested here \\
\midrule
C1 & SIPIC effective dielectric \& down-shift (Eqs.~12/16) & analytic + numerical & Yes & \textbf{Yes (core)} \\
C2 & Modified Bohm--Gross $\omega^{2}=(\omega_{pe,\text{SI}})^{2}+1.5\,v_{th}^{2}k^{2}$ (Eq.~18) & analytic + numerical & Yes & Partial (cold limit of C1) \\
C3 & Stable $\forall\,\Delta t$ if $C_{\text{SI}}\ge 1$ ($\omega\Delta t < 2/\sqrt{C_{\text{SI}}}$, Eq.~14) & stability & Yes & Yes (implicit) \\
C4 & Upper/lower hybrid frequencies unaffected except via $\omega_{pe,\text{SI}}$ (Eqs.~19--20) & numerical & Yes (2D B-field) & No \\
C5 & Total energy change $<2.5\%$; SIPIC$\approx$ECPIC & numerical & Yes & No (needs WarpX/Aleph) \\
C6 & Landau damping rates preserved (\S III.C) & numerical & Yes & No \\
\bottomrule
\end{tabular}
\end{center}

\noindent\textbf{Focus: C1} (central verification claim), with C3 as a
byproduct.

\section{Method}
Independent re-implementation (no paper code; a $\sim$180-line NumPy PIC
in \texttt{work/sipic\_dispersion.py}):
\begin{enumerate}
  \item 1D electrostatic PIC, periodic box, normalized units
        ($\omega_{pe}=1$, $\varepsilon_0=1$, $q=-1$, $m=1$).
        Leapfrog particle push; Cloud-In-Cell (linear) charge deposition
        and field gather; spectral (FFT) Poisson solve.
  \item SIPIC modification implemented exactly as the paper's field
        operator (Eqs.~9--10): the negative-Laplacian is multiplied by
        $F = (1 + C_{\text{SI}}\,\omega_{pe}^{2}\,\Delta t^{2}/4)$,
        i.e.\ $\varepsilon_{\text{eff}} = \varepsilon_0\,F$. The paper's
        Eq.~12 writes $\kappa = 1/F$, $\varepsilon_{\text{SI}} =
        \varepsilon_0\kappa$, $\omega^{2} = \omega_{p}^{2}/\kappa$; that
        $\kappa$-bookkeeping is internally inverted relative to the
        \emph{stated} Eq.~16 down-shift, so we implement the physical
        operator from Eqs.~9--10 directly. The resulting down-shift
        $1/\sqrt{F}$ is identical to Eq.~16.
  \item Test problem: cold Langmuir oscillation --- uniform plasma with
        a small mode-1 position perturbation (\texttt{amp} = 0.02);
        $v_{th}=0$ (cold, so the Bohm--Gross thermal term vanishes and
        C1 is isolated); $N=80{,}000$ particles, 32 cells, long box
        (small $k$).
  \item Sweep: numerical factor $a\in\{0.5,1,2,4,8\}$ giving
        $\omega_{pe}\,\Delta t\in\{1,2,4,8,16\}$ (16$\times$ range
        spanning both $\Delta t$ conventions in the paper), with
        $C_{\text{SI}}=4$ (Table~I).
  \item Diagnostic: record the signed real part of the complex mode-1
        Fourier coefficient of $E$ each step; FFT the time series (Hann
        window, parabolic sub-bin interpolation, physical search band)
        to extract $\omega$.
  \item Validation control: run classical PIC ($C_{\text{SI}}=0$) at
        well-resolved $\Delta t$ --- a correct code must recover
        $\omega\approx\omega_{pe}$.
  \item Judge: free Argo \texttt{argo:gpt-5.2} (localhost:44497),
        temperature 0, given the measured-vs-analytic table (no regex
        scoring).
\end{enumerate}
Tools: Python\,3 + NumPy + Matplotlib (local, CherryRd). OA PDF fetch +
\texttt{pdftotext} via \texttt{ssh uicgpu}. Free endpoints only.

\section{Results vs paper}

\subsection*{Validation (classical PIC, $C_{\text{SI}}=0$) --- sanity check}
\begin{center}
\begin{tabular}{@{}ccc@{}}
\toprule
$\omega_{pe}\,\Delta t$ & measured $\omega/\omega_{pe}$ & target \\
\midrule
0.10 & 1.033 & 1.0 \\
0.20 & 1.011 & 1.0 \\
0.50 & 1.017 & 1.0 \\
\bottomrule
\end{tabular}
\end{center}
The from-scratch PIC recovers the physical plasma frequency to
$\sim$1--3\%. Diagnostic sound.

\subsection*{C1 --- SIPIC down-shift ($C_{\text{SI}}=4$): measured vs Eq.~16}
\begin{center}
\begin{tabular}{@{}ccccc@{}}
\toprule
$\omega_{pe}\,\Delta t$ & measured $\omega/\omega_{pe}$ & Eq.~16 prediction & classical (unmodified) & error vs Eq.~16 \\
\midrule
1  & 0.731 & 0.707 & 1.000 & 3.3\,\% \\
2  & 0.451 & 0.447 & 1.000 & 0.9\,\% \\
4  & 0.252 & 0.243 & 1.000 & 3.9\,\% \\
8  & 0.126 & 0.124 & 1.000 & 1.9\,\% \\
16 & 0.069 & 0.062 & 1.000 & 10.2\,\% \\
\bottomrule
\end{tabular}
\end{center}

The measured plasma-oscillation frequency tracks the analytic SIPIC
down-shift (Eq.~16) across a 16$\times$ range of $\omega_{pe}\,\Delta t$
--- $\sim$1--4\% for $\omega_{pe}\,\Delta t\le 8$ and $\sim$10\% at the
most extreme step ($\omega_{pe}\,\Delta t = 16$, only $\sim$6 samples
per down-shifted oscillation) --- and is decisively inconsistent with
the classical, unmodified $\omega/\omega_{pe}=1$. This is exactly the
paper's stated verification result (``As $\Delta t$ is increased, the
contours follow the modified dispersion relation \ldots the modified
dispersion curves traverse downward,'' \S III.A). Plot:
\texttt{evidence/sipic\_downshift.png}. Raw numbers:
\texttt{evidence/sipic\_dispersion\_results.json}.

\subsection*{Stability (C3)}
All runs, including $\omega_{pe}\,\Delta t=16$ (16$\times$ beyond the
explicit-PIC limit), remained bounded with the mode amplitude
oscillating cleanly --- consistent with the paper's stability bound
$\omega\Delta t < 2/\sqrt{C_{\text{SI}}}$ for $C_{\text{SI}}\ge 1$
(Eq.~14).

\subsection*{Discrepancies / caveats}
\begin{itemize}
  \item \textbf{Paper-internal inconsistency:} Table~I gives $\Delta t =
        2a/\omega_{pe}$ ($\omega_{pe}\Delta t = 2a$) while \S III.A
        prose gives $\Delta t\cdot\omega_{pe} = a/2$. Because Eq.~16
        depends only on the product $\omega_{pe}\,\Delta t$, the
        replication sweeps the full range and the conclusion is
        independent of this ambiguity.
  \item \textbf{Paper $\kappa$-notation inversion:} Eq.~12/13
        ($\kappa=1/F$, $\omega^{2}=\omega_{p}^{2}/\kappa$) written
        literally gives an \emph{up}-shift, contradicting the Eq.~16
        text (``reduced frequency'') and Fig.~1 (downward contours).
        The physically correct + self-consistent reading (implemented
        here) is the down-shift; this is a typographic/bookkeeping
        subtlety, not a physics error.
  \item Not tested: hybrid modes (C4), full energy-conservation curves
        (C5, needs WarpX/Aleph), Landau damping (C6). C2's thermal
        Bohm--Gross term was intentionally suppressed (cold run) to
        isolate C1.
\end{itemize}

\section{LLM-judge verdict (free Argo gpt-5.2, verbatim)}
\begin{quote}
``The measured $\omega/\omega_{pe}$ agrees with Eq.~16 within
$\sim$1--4\% for $\omega_{pe}\,\Delta t\le 8$ and shows the correct
down-shift trend at 16 (but with $\sim$10\% error) \ldots the core
claim --- an effective dielectric producing a stable, down-shifted
plasma frequency scaling like Eq.~16 for large $\omega_{pe}\,\Delta t$
--- is reproduced. \textbf{VERDICT: REPLICATED.}''
\end{quote}
Full text: \texttt{evidence/llm\_judge\_verdict.txt}.

\section{GENUINE CRITIQUE}
This section deliberately steps outside the celebratory framing to
enumerate real weaknesses of \emph{this} replication effort. Nothing
here overturns the verdict, but each point identifies a way the current
evidence is thinner than the executive summary might suggest.

\paragraph{1.\ Only one of six claims was actually re-executed.}
C1 is the central verification claim, but the paper's argument spans
six claims (C1--C6) covering hybrid modes, energy conservation, and
Landau damping. This replication tested C1 quantitatively and inherited
C3 (stability) as a byproduct of the same runs. C2 was reduced to its
cold limit (thermal term suppressed), and C4--C6 were not attempted at
all. A ``REPLICATED'' verdict at the paper level therefore rests on
one-and-a-fraction of six claims. The verdict is honestly scoped to C1
in the body of the report, but any downstream reuse of the verdict
should carry that scope with it.

\paragraph{2.\ Cold-plasma, 1D, small-mode-1 setup is the friendliest
possible test.} $v_{th}=0$ zeros the Bohm--Gross thermal term, and a
mode-1 perturbation with $N=80{,}000$ CIC particles has an unusually
clean signal-to-noise ratio compared to a warm, multi-mode, or 2D/3D
run. This is appropriate for isolating C1, but it means the
replication does not stress the SIPIC operator in the regimes where a
production plasma code would actually live (kinetic tails, magnetized
sheaths, multi-species). Agreement in this friendly limit is necessary
but not sufficient evidence that the paper's scheme behaves as
advertised in realistic use.

\paragraph{3.\ The 10\,\% error at $\omega_{pe}\Delta t = 16$ is
under-characterized.} With only $\sim$6 samples per down-shifted
oscillation at that step, the FFT-plus-parabolic-interpolation
diagnostic is near its resolution floor. The 10\,\% number is reported
as an error against Eq.~16, but it could reflect (a) a real
scheme-vs-theory departure at extreme $\Delta t$, (b) a diagnostic
artifact (windowing/aliasing/sub-bin interpolation bias), or (c) both.
The replication does not distinguish these, and does not report
uncertainty bars from repeated runs with different seeds/particle
counts. That would be the natural next test.

\paragraph{4.\ The $\kappa$-notation resolution is a judgement call.}
The paper's Eq.~12/13 literally implies an \emph{up}-shift, which
contradicts Eq.~16 and Fig.~1. The replication assumes this is a
typo/bookkeeping inversion and implements the physical operator from
Eqs.~9--10, which produces the observed down-shift. This is almost
certainly the intended reading --- the numerical agreement supports it
--- but a properly rigorous replication would either flag this back to
the authors for confirmation or explicitly test both readings and
report which one the paper's own Fig.~1 contours are consistent with.

\paragraph{5.\ No cross-code comparison; original codes not touched.}
The paper's verification is partly a WarpX-vs-Aleph code-to-code
exercise. This replication is a from-scratch NumPy PIC that touches
neither WarpX nor Aleph. That is deliberate (from-scratch avoids
inheriting bugs from either code), but it also means we cannot
adjudicate the paper's cross-code agreement claims --- only the
underlying analytic prediction. If WarpX and Aleph disagreed with each
other on some detail, this replication would not have seen it.

\paragraph{6.\ Judge is a single LLM at $T=0$, not a panel.}
Using free Argo \texttt{argo:gpt-5.2} at temperature 0 with the
measured-vs-analytic table in hand is a reasonable last-mile check, but
it is one judge, one prompt, one run. The verdict is not blind (the
judge sees the table already framed as ``measured'' vs ``prediction''),
and there is no adversarial second judge or human physicist in the
loop. The judge's role here is closer to a sanity check on the
arithmetic than an independent verdict.

\paragraph{7.\ ``REPLICATED'' should be read as ``C1 reproduced in the
cold-limit 1D setting to $\sim$1--4\% for $\omega_{pe}\Delta t\le 8$.''}
That is what the numbers actually support. Stronger readings
(``the paper's scheme works,'' ``the paper's verification is
independently confirmed'') go beyond what was tested and should be
avoided.

\section{Verdict}
\textbf{Verdict: REPLICATED (scoped to C1 in the cold 1D limit).}

The paper's central verification claim --- that the SIPIC
effective-dielectric modification down-shifts the electron plasma
frequency by $1/\sqrt{1 + C_{\text{SI}}\,\omega_{pe}^{2}\,\Delta t^{2}/4}$
(Eqs.~12/16), enabling stable operation far beyond the explicit
$\omega_{pe}\,\Delta t$ limit --- was independently reproduced from the
paper's equations with a from-scratch 1D ES-PIC. Measured frequencies
track the analytic prediction to $\sim$1--4\% over
$\omega_{pe}\,\Delta t\le 8$ (10\,\% at 16) and are decisively distinct
from the unmodified classical $\omega_{pe}$; a classical-PIC control
recovers $\omega\approx\omega_{pe}$.

\bigskip
\noindent\textbf{WAVE\_RESULT} set=OSTI-100 paper=3374709
verdict=REPLICATED
dir=\~/Dropbox/REPLICATE-PROJECT/OSTI-3374709-semi-implicit-ecpic-verification
one\_line=Independent 1D ES-PIC reproduces the SIPIC plasma-frequency
down-shift $\omega_{pe}/\sqrt{1+C_{\text{SI}}\omega_{pe}^{2}\Delta t^{2}/4}$
(Eqs.~12/16) to $\sim$1--4\% over $\omega_{pe}\,\Delta t\le 8$
(10\,\% at 16), decisively unlike classical $\omega_{pe}$.

\end{document}
